% Options for packages loaded elsewhere
% Options for packages loaded elsewhere
\PassOptionsToPackage{unicode}{hyperref}
\PassOptionsToPackage{hyphens}{url}
\PassOptionsToPackage{dvipsnames,svgnames,x11names}{xcolor}
%
\documentclass[
  letterpaper,
  DIV=11,
  numbers=noendperiod]{scrartcl}
\usepackage{xcolor}
\usepackage{amsmath,amssymb}
\setcounter{secnumdepth}{-\maxdimen} % remove section numbering
\usepackage{iftex}
\ifPDFTeX
  \usepackage[T1]{fontenc}
  \usepackage[utf8]{inputenc}
  \usepackage{textcomp} % provide euro and other symbols
\else % if luatex or xetex
  \usepackage{unicode-math} % this also loads fontspec
  \defaultfontfeatures{Scale=MatchLowercase}
  \defaultfontfeatures[\rmfamily]{Ligatures=TeX,Scale=1}
\fi
\usepackage{lmodern}
\ifPDFTeX\else
  % xetex/luatex font selection
\fi
% Use upquote if available, for straight quotes in verbatim environments
\IfFileExists{upquote.sty}{\usepackage{upquote}}{}
\IfFileExists{microtype.sty}{% use microtype if available
  \usepackage[]{microtype}
  \UseMicrotypeSet[protrusion]{basicmath} % disable protrusion for tt fonts
}{}
\makeatletter
\@ifundefined{KOMAClassName}{% if non-KOMA class
  \IfFileExists{parskip.sty}{%
    \usepackage{parskip}
  }{% else
    \setlength{\parindent}{0pt}
    \setlength{\parskip}{6pt plus 2pt minus 1pt}}
}{% if KOMA class
  \KOMAoptions{parskip=half}}
\makeatother
% Make \paragraph and \subparagraph free-standing
\makeatletter
\ifx\paragraph\undefined\else
  \let\oldparagraph\paragraph
  \renewcommand{\paragraph}{
    \@ifstar
      \xxxParagraphStar
      \xxxParagraphNoStar
  }
  \newcommand{\xxxParagraphStar}[1]{\oldparagraph*{#1}\mbox{}}
  \newcommand{\xxxParagraphNoStar}[1]{\oldparagraph{#1}\mbox{}}
\fi
\ifx\subparagraph\undefined\else
  \let\oldsubparagraph\subparagraph
  \renewcommand{\subparagraph}{
    \@ifstar
      \xxxSubParagraphStar
      \xxxSubParagraphNoStar
  }
  \newcommand{\xxxSubParagraphStar}[1]{\oldsubparagraph*{#1}\mbox{}}
  \newcommand{\xxxSubParagraphNoStar}[1]{\oldsubparagraph{#1}\mbox{}}
\fi
\makeatother

\usepackage{color}
\usepackage{fancyvrb}
\newcommand{\VerbBar}{|}
\newcommand{\VERB}{\Verb[commandchars=\\\{\}]}
\DefineVerbatimEnvironment{Highlighting}{Verbatim}{commandchars=\\\{\}}
% Add ',fontsize=\small' for more characters per line
\usepackage{framed}
\definecolor{shadecolor}{RGB}{241,243,245}
\newenvironment{Shaded}{\begin{snugshade}}{\end{snugshade}}
\newcommand{\AlertTok}[1]{\textcolor[rgb]{0.68,0.00,0.00}{#1}}
\newcommand{\AnnotationTok}[1]{\textcolor[rgb]{0.37,0.37,0.37}{#1}}
\newcommand{\AttributeTok}[1]{\textcolor[rgb]{0.40,0.45,0.13}{#1}}
\newcommand{\BaseNTok}[1]{\textcolor[rgb]{0.68,0.00,0.00}{#1}}
\newcommand{\BuiltInTok}[1]{\textcolor[rgb]{0.00,0.23,0.31}{#1}}
\newcommand{\CharTok}[1]{\textcolor[rgb]{0.13,0.47,0.30}{#1}}
\newcommand{\CommentTok}[1]{\textcolor[rgb]{0.37,0.37,0.37}{#1}}
\newcommand{\CommentVarTok}[1]{\textcolor[rgb]{0.37,0.37,0.37}{\textit{#1}}}
\newcommand{\ConstantTok}[1]{\textcolor[rgb]{0.56,0.35,0.01}{#1}}
\newcommand{\ControlFlowTok}[1]{\textcolor[rgb]{0.00,0.23,0.31}{\textbf{#1}}}
\newcommand{\DataTypeTok}[1]{\textcolor[rgb]{0.68,0.00,0.00}{#1}}
\newcommand{\DecValTok}[1]{\textcolor[rgb]{0.68,0.00,0.00}{#1}}
\newcommand{\DocumentationTok}[1]{\textcolor[rgb]{0.37,0.37,0.37}{\textit{#1}}}
\newcommand{\ErrorTok}[1]{\textcolor[rgb]{0.68,0.00,0.00}{#1}}
\newcommand{\ExtensionTok}[1]{\textcolor[rgb]{0.00,0.23,0.31}{#1}}
\newcommand{\FloatTok}[1]{\textcolor[rgb]{0.68,0.00,0.00}{#1}}
\newcommand{\FunctionTok}[1]{\textcolor[rgb]{0.28,0.35,0.67}{#1}}
\newcommand{\ImportTok}[1]{\textcolor[rgb]{0.00,0.46,0.62}{#1}}
\newcommand{\InformationTok}[1]{\textcolor[rgb]{0.37,0.37,0.37}{#1}}
\newcommand{\KeywordTok}[1]{\textcolor[rgb]{0.00,0.23,0.31}{\textbf{#1}}}
\newcommand{\NormalTok}[1]{\textcolor[rgb]{0.00,0.23,0.31}{#1}}
\newcommand{\OperatorTok}[1]{\textcolor[rgb]{0.37,0.37,0.37}{#1}}
\newcommand{\OtherTok}[1]{\textcolor[rgb]{0.00,0.23,0.31}{#1}}
\newcommand{\PreprocessorTok}[1]{\textcolor[rgb]{0.68,0.00,0.00}{#1}}
\newcommand{\RegionMarkerTok}[1]{\textcolor[rgb]{0.00,0.23,0.31}{#1}}
\newcommand{\SpecialCharTok}[1]{\textcolor[rgb]{0.37,0.37,0.37}{#1}}
\newcommand{\SpecialStringTok}[1]{\textcolor[rgb]{0.13,0.47,0.30}{#1}}
\newcommand{\StringTok}[1]{\textcolor[rgb]{0.13,0.47,0.30}{#1}}
\newcommand{\VariableTok}[1]{\textcolor[rgb]{0.07,0.07,0.07}{#1}}
\newcommand{\VerbatimStringTok}[1]{\textcolor[rgb]{0.13,0.47,0.30}{#1}}
\newcommand{\WarningTok}[1]{\textcolor[rgb]{0.37,0.37,0.37}{\textit{#1}}}

\usepackage{longtable,booktabs,array}
\usepackage{calc} % for calculating minipage widths
% Correct order of tables after \paragraph or \subparagraph
\usepackage{etoolbox}
\makeatletter
\patchcmd\longtable{\par}{\if@noskipsec\mbox{}\fi\par}{}{}
\makeatother
% Allow footnotes in longtable head/foot
\IfFileExists{footnotehyper.sty}{\usepackage{footnotehyper}}{\usepackage{footnote}}
\makesavenoteenv{longtable}
\usepackage{graphicx}
\makeatletter
\newsavebox\pandoc@box
\newcommand*\pandocbounded[1]{% scales image to fit in text height/width
  \sbox\pandoc@box{#1}%
  \Gscale@div\@tempa{\textheight}{\dimexpr\ht\pandoc@box+\dp\pandoc@box\relax}%
  \Gscale@div\@tempb{\linewidth}{\wd\pandoc@box}%
  \ifdim\@tempb\p@<\@tempa\p@\let\@tempa\@tempb\fi% select the smaller of both
  \ifdim\@tempa\p@<\p@\scalebox{\@tempa}{\usebox\pandoc@box}%
  \else\usebox{\pandoc@box}%
  \fi%
}
% Set default figure placement to htbp
\def\fps@figure{htbp}
\makeatother





\setlength{\emergencystretch}{3em} % prevent overfull lines

\providecommand{\tightlist}{%
  \setlength{\itemsep}{0pt}\setlength{\parskip}{0pt}}



 


\usepackage{makeidx}
\makeindex
\KOMAoption{captions}{tableheading}
\makeatletter
\@ifpackageloaded{caption}{}{\usepackage{caption}}
\AtBeginDocument{%
\ifdefined\contentsname
  \renewcommand*\contentsname{Table of contents}
\else
  \newcommand\contentsname{Table of contents}
\fi
\ifdefined\listfigurename
  \renewcommand*\listfigurename{List of Figures}
\else
  \newcommand\listfigurename{List of Figures}
\fi
\ifdefined\listtablename
  \renewcommand*\listtablename{List of Tables}
\else
  \newcommand\listtablename{List of Tables}
\fi
\ifdefined\figurename
  \renewcommand*\figurename{Figure}
\else
  \newcommand\figurename{Figure}
\fi
\ifdefined\tablename
  \renewcommand*\tablename{Table}
\else
  \newcommand\tablename{Table}
\fi
}
\@ifpackageloaded{float}{}{\usepackage{float}}
\floatstyle{ruled}
\@ifundefined{c@chapter}{\newfloat{codelisting}{h}{lop}}{\newfloat{codelisting}{h}{lop}[chapter]}
\floatname{codelisting}{Listing}
\newcommand*\listoflistings{\listof{codelisting}{List of Listings}}
\makeatother
\makeatletter
\makeatother
\makeatletter
\@ifpackageloaded{caption}{}{\usepackage{caption}}
\@ifpackageloaded{subcaption}{}{\usepackage{subcaption}}
\makeatother
\usepackage{bookmark}
\IfFileExists{xurl.sty}{\usepackage{xurl}}{} % add URL line breaks if available
\urlstyle{same}
\hypersetup{
  pdftitle={Minimal Repro: Multi-line LaTeX macro + footnote → Extra \}},
  colorlinks=true,
  linkcolor={blue},
  filecolor={Maroon},
  citecolor={Blue},
  urlcolor={Blue},
  pdfcreator={LaTeX via pandoc}}


\title{Minimal Repro: Multi-line LaTeX macro + footnote → Extra \}}
\usepackage{etoolbox}
\makeatletter
\providecommand{\subtitle}[1]{% add subtitle to \maketitle
  \apptocmd{\@title}{\par {\large #1 \par}}{}{}
}
\makeatother
\subtitle{Tests whether parse-latex expands user-defined macros during
PDF builds, and whether pandoc itself expands them independently}
\author{}
\date{2026-04-14}
\begin{document}
\maketitle


\subsection{Context}\label{context}

In a Quarto book (\texttt{HTML\ +\ PDF}), every chapter includes
\texttt{latex/latex-commands.qmd} via
\texttt{\{\{\textless{}\ include\ \textgreater{}\}\}}. That file defines
fallback index macros for HTML contexts:

\begin{Shaded}
\begin{Highlighting}[]
\FunctionTok{\textbackslash{}providecommand}\NormalTok{\{}\ExtensionTok{\textbackslash{}ixd}\NormalTok{\}[1]\{}\CommentTok{\%}
  \FunctionTok{\textbackslash{}index}\NormalTok{\{\#1@}\FunctionTok{\textbackslash{}texttt}\NormalTok{\{\#1\} data\}}\CommentTok{\%}
  \FunctionTok{\textbackslash{}index}\NormalTok{\{datasets!\#1@}\FunctionTok{\textbackslash{}texttt}\NormalTok{\{\#1\}\}}\CommentTok{\%}
\NormalTok{\}}
\end{Highlighting}
\end{Shaded}

Standard LaTeX multi-line style: \texttt{\%} at line-ends suppresses
whitespace in the macro \emph{expansion} (XeLaTeX interprets \texttt{\%}
as a comment, discards the newline). This is correct and normal for
LaTeX. When XeLaTeX expands \texttt{\textbackslash{}ixd\{Prestige\}}, it
produces
\texttt{\textbackslash{}index\{...\}\textbackslash{}index\{...\}} with
no \texttt{\%} in the result.

\textbf{The problem:} Something during the \emph{pandoc pass} expands
\texttt{\textbackslash{}ixd\{Prestige\}} \emph{literally}, inserting the
\texttt{\%} characters from the definition body into the output text.
When a Markdown footnote (\texttt{.{[}\^{}fn{]}}) follows in the source,
the output becomes:

\begin{Shaded}
\begin{Highlighting}[]
\FunctionTok{\textbackslash{}index}\NormalTok{\{Prestige@}\FunctionTok{\textbackslash{}texttt}\NormalTok{\{Prestige\} data\}}\CommentTok{\%}
\FunctionTok{\textbackslash{}index}\NormalTok{\{datasets!Prestige@}\FunctionTok{\textbackslash{}texttt}\NormalTok{\{Prestige\}\}}\CommentTok{\%.\textbackslash{}footnote\{...\}}
\end{Highlighting}
\end{Shaded}

The \texttt{\%.\textbackslash{}footnote\{} means XeLaTeX never opens the
footnote →
\texttt{!\ Extra\ \},\ or\ forgotten\ \textbackslash{}endgroup}.

\textbf{Setup:} \texttt{parse-latex} is in
\texttt{format:\ html:\ filters:} only (not global \texttt{filters:}).
Per @mcanouil, \texttt{FORMAT\ =\ "latex"} during the PDF pandoc pass,
so the filter's
\texttt{if\ not\ FORMAT:match\ \textquotesingle{}html\textquotesingle{}\ then\ return\ \{\}\ end}
guard should fire and do nothing. \textbf{If that is correct, the
expansion must come from a different mechanism.}

\begin{center}\rule{0.5\linewidth}{0.5pt}\end{center}

\subsection{Test 1: Static macro + footnote (parse-latex in HTML filters
only)}\label{sec-test1}

A raw LaTeX block defines a multi-line macro; a raw inline uses it; a
Markdown footnote follows immediately. This is the minimal reproduction
of the book scenario.

\providecommand{\myone}[1]{%
  \textit{T1: #1}%
}

Static inline then footnote: \myone{alpha}.\footnote{Test 1 footnote.
  \textbf{If PDF build fails here} with \texttt{Extra\ \}}, something
  expanded \texttt{\textbackslash{}myone} during the pandoc pass
  (inserting literal \texttt{\%}). Check the generated \texttt{.tex}
  file for: - PASS:
  \texttt{\textbackslash{}myone\{alpha\}.\textbackslash{}footnote\{Test\ 1\ footnote...\}}
  (unexpanded, one line) - FAIL:
  \texttt{\textbackslash{}textit\{T1:\ alpha\}\%\textbackslash{}n.\textbackslash{}footnote\{...\}}
  (expanded, \texttt{\%} before footnote)}

\begin{center}\rule{0.5\linewidth}{0.5pt}\end{center}

\subsection{\texorpdfstring{Test 2: Definition via
\texttt{\{\{\textless{}\ include\ \textgreater{}\}\}}}{Test 2: Definition via \{\{\textless{} include \textgreater\}\}}}\label{sec-test2}

In the actual book, the \texttt{\textbackslash{}providecommand} comes
from an included file, not directly in the chapter source. Quarto
resolves \texttt{\{\{\textless{}\ include\ \textgreater{}\}\}} before
the pandoc pass, so the included content ends up in the freeze cache
(\texttt{tex.json}). Any stale cache will embed the \emph{old}
multi-line definition even after the source file is updated.

\providecommand{\mytwo}[1]{%
  \textit{T2: #1}%
}
\providecommand{\mythree}[1]{%
  \textit{T3: #1}%
}

\mytwo{beta}.\footnote{Test 2 footnote. If this fails but Test 1 passes,
  the issue is specific to
  \texttt{\{\{\textless{}\ include\ \textgreater{}\}\}}'d definitions
  (e.g., a stale freeze cache embedding the definition from a previous
  build).}

\begin{center}\rule{0.5\linewidth}{0.5pt}\end{center}

\subsection{Test 3: Inline R output near footnote}\label{sec-test3}

The book uses inline R expressions that emit raw LaTeX macros for PDF:

\begin{Shaded}
\begin{Highlighting}[]
\FunctionTok{dataset}\NormalTok{(}\StringTok{"Prestige"}\NormalTok{)  →  }\StringTok{"}\SpecialCharTok{\textbackslash{}t}\StringTok{exttt\{Prestige\}\textbackslash{}ixd\{Prestige\}"}\NormalTok{  (}\ControlFlowTok{for}\NormalTok{ PDF)}
\end{Highlighting}
\end{Shaded}

This test simulates that without the full book infrastructure:

\mythree{gamma}

.\footnote{Test 3 footnote. If this fails but Test 1 passes, the issue
  is specific to knitr-emitted raw LaTeX content and how Quarto's freeze
  cache stores it. (\texttt{\textbackslash{}mythree} is defined in
  \texttt{quarto-pandoc-test-macros.qmd}, included in Test 2.)}

\begin{center}\rule{0.5\linewidth}{0.5pt}\end{center}

\subsection{How to run}\label{how-to-run}

\begin{Shaded}
\begin{Highlighting}[]
\CommentTok{\# From the project root:}
\ExtensionTok{quarto}\NormalTok{ render issues/quarto{-}pandoc{-}test.qmd }\AttributeTok{{-}{-}to}\NormalTok{ pdf}
\ExtensionTok{quarto}\NormalTok{ render issues/quarto{-}pandoc{-}test.qmd }\AttributeTok{{-}{-}to}\NormalTok{ html}
\end{Highlighting}
\end{Shaded}

\textbf{To test WITHOUT parse-latex} (isolates pandoc's own behavior):
Remove \texttt{filters:\ {[}tarleb/parse-latex{]}} from the YAML and
rebuild. If the error persists, pandoc itself is expanding the macros
--- parse-latex is not involved.

\textbf{To test WITHOUT
\texttt{\{\{\textless{}\ include\ \textgreater{}\}\}}}: Remove the
\texttt{\{\{\textless{}\ include\ \textgreater{}\}\}} line and
\texttt{\textbackslash{}mytwo}/\texttt{\textbackslash{}mythree} tests.

\textbf{To diagnose which mechanism is expanding macros}, add this Lua
filter alongside parse-latex (per @mcanouil's suggestion) to confirm
FORMAT during the PDF pass:

\begin{Shaded}
\begin{Highlighting}[]
\CommentTok{{-}{-} debug{-}format.lua}
\KeywordTok{function}\NormalTok{ Pandoc}\OperatorTok{(}\VariableTok{doc}\OperatorTok{)}
  \FunctionTok{io.stderr}\OperatorTok{:}\FunctionTok{write}\OperatorTok{(}\StringTok{"[debug] FORMAT = "} \OperatorTok{..} \FunctionTok{tostring}\OperatorTok{(}\ConstantTok{FORMAT}\OperatorTok{)} \OperatorTok{..} \StringTok{"}\SpecialCharTok{\textbackslash{}n}\StringTok{"}\OperatorTok{)}
  \ControlFlowTok{return} \VariableTok{doc}
\KeywordTok{end}
\end{Highlighting}
\end{Shaded}

\begin{center}\rule{0.5\linewidth}{0.5pt}\end{center}

\subsection{\texorpdfstring{Expected \texttt{.tex} content (with
\texttt{keep-tex:\ true})}{Expected .tex content (with keep-tex: true)}}\label{expected-.tex-content-with-keep-tex-true}

\textbf{If macros are NOT expanded during pandoc pass} (correct):

\begin{Shaded}
\begin{Highlighting}[]
\FunctionTok{\textbackslash{}myone}\NormalTok{\{alpha\}.}\FunctionTok{\textbackslash{}footnote}\NormalTok{\{Test 1 footnote...\}}
\FunctionTok{\textbackslash{}mytwo}\NormalTok{\{beta\}.}\FunctionTok{\textbackslash{}footnote}\NormalTok{\{Test 2 footnote...\}}
\end{Highlighting}
\end{Shaded}

XeLaTeX then expands these natively; the \texttt{\%} in the definition
body acts as a TeX comment (discards the newline), so the expansion
contains no \texttt{\%}.

\textbf{If macros ARE expanded during pandoc pass} (the bug):

\begin{Shaded}
\begin{Highlighting}[]
\FunctionTok{\textbackslash{}textit}\NormalTok{\{T1: alpha\}}\CommentTok{\%}
\NormalTok{.}\FunctionTok{\textbackslash{}footnote}\NormalTok{\{Test 1 footnote...\}}
\end{Highlighting}
\end{Shaded}

The literal \texttt{\%} from the definition body is inserted into the
pandoc output, commenting out the \texttt{.\textbackslash{}footnote\{}
opener.




\end{document}
